\documentclass[../main.tex]{subfiles}
\begin{document}
%==========================================TIÊU ĐỀ TẬP========================================
\begin{table}[h]
  \begin{adjustwidth}{-1cm}{}
    \begin{tabular}{>{\centering\arraybackslash}p{6cm}|>{\raggedright\arraybackslash}p{12.5cm}}
      \multirow{5}{6cm}
      % Cột bên trái: ÉPISODE và số tập
      {\\[-40pt]
        \flaregothic\fontsize{20pt}{20pt}\selectfont \centering
        \color{eptype}ÉPISODE\color{black}\\
        \burbank\fontsize{72pt}{72pt}\selectfont
        \color{epnum}34\color{black}\\[6pt]
        \cafeteria\fontsize{20pt}{20pt}\selectfont\color{epnum}(3-5)\color{black}
        \\[18pt]
        \flaregothic\fontsize{18pt}{18pt}\selectfont
        \color{epattr}{\color{comiket}{Comiket d'hiver}}\\
        \flaregothic\fontsize{14pt}{14pt}\selectfont
        \color{epattr}BIENVENUE, \uline{\color{Pekopeko} Pekora} !\\
        \flaregothic\fontsize{14pt}{14pt}\selectfont
        \color{epattr}LE DERNIER ÉPISODE DE L'ANNÉE 2019 !\\[24pt]
        \color{black}
      }
       &
      % Dòng thứ nhất: tên tập tiếng Nhật
      {\fotskip\fontsize{18pt}{18pt}\selectfont \ruby{壁}{かべ}に\ruby{埋}{う}まりました。\ruby{助}{たす}けてください
      }                           \\[6pt]
      % Dòng thứ hai: tên tập tiếng Anh
       & {\cafeteria\fontsize{18pt}{18pt}\selectfont Buried in Wall. Get Some Help! (Winter Comiket)}                     \\[4pt]
      % Dòng thứ ba: tên tập tiếng Việt
       & {\vnmsans\fontsize{18pt}{18pt}\selectfont Kịch Comiket II - Shion bị kẹt tường}                                  \\[8pt]
      % Dòng thứ tư: Nhân vật xuất hiện
       & {\caviar{\fontsize{14pt}{14pt}\selectfont Nhân vật: \emp{kuon1} \emp{shion} \emp{korone} \emp{pekora} \emp{noel}}
         }
      \\[8pt]
      % Dòng cuối cùng: Thời gian diễn ra của tập (thường sẽ là ngày phát hành)
       & {\continuum\fontsize{16pt}{16pt}\selectfont Ngày phát hành: 29/12/2019}                                          \\[18pt]
    \end{tabular}
  \end{adjustwidth}
\end{table}
%=========================================TRUYỆN CHÍNH========================================
\color{black}

\txtbx[2]{{\color{vol} \uline{Tóm tắt:}} Sắp đến Comiket mùa đông, nhưng không khí bỗng chốc trở nên hỗn loạn khi Shion bị kẹt trong bức tường. Sự tò điên về củ cà rốt của Pekora đã vô tình mở ra một hố đen, hút mọi người vào một thế giới kỳ dị bên trong\dots\ lọn tóc của nàng Thỏ. Một cuộc giải cứu dở khóc dở cười diễn ra, kết thúc bằng một vụ kẹt tường khác.}

Chỉ còn đúng hai ngày nữa là khép lại năm 2019. Không khí cuối năm rộn ràng, tất bật bao trùm khắp nơi, và \hll\ đang rục rịch chuẩn bị cho Comiket mùa đông - sự kiện được mong chờ nhất trong năm. Trong phòng nghỉ, Korone, Noel, Kuon và một thành viên mới của Thế hệ thứ Ba đang tụ tập bàn luận rôm rả.

\chrint

\model{25}{l}{6cm}{pekora}

Usada Pekora (\ruby{兎}{うさ}\ruby{田}{だ}ぺこら, \ruby{U-xa}{Thỏ}-\ruby{đa}{Điền}\ Pê-cô-ra) là một cô thỏ đến từ Pekoland, nơi cô luôn tự nhận mình là thành viên của hoàng gia (ít nhất, theo như CV của cô ghi lại). Cô là một người đơn độc, cực kỳ, CỰC KỲ yêu cà rốt\footnote{Mặc dù về thực tế, cà rốt không phải là thức ăn sở trường của thỏ, và chúng không bao giờ ăn nhiều đến mức đó.}, tới mức dù đi đâu cô cũng phải giắt theo vài củ bên mình.

Tính cách của Pekora có thể được mô tả là bề ngoài thì trẻ con, kiêu ngạo, non nớt và khó tính, nhưng sâu thẳm lại thực sự vui tươi và thân thiện. Khi không lên hình, cô được những thành viên \hll\ khác nhận xét là rất lịch sự và nhút nhát, cũng như là người nội tâm hóa nỗi đau của mình - điều từng khiến một trong những buổi phát trực tiếp của cô phải kết thúc sớm với cảnh cô òa khóc nức nở.

Pekora là một cô nàng vô cùng tinh quái, đến mức các thành viên khác trong công ty (và cả Kuon) sẽ lập tức nghi ngờ bất kỳ điều gì mà cô nhúng tay vào. Mặc dù Miko thường là nạn nhân số một của những trò lừa gạt do cô bày ra, nhưng ngay cả bản thân Pekora đôi khi cũng không thể thoát khỏi gậy ông đập lưng ông. Dù vậy, Pekora vẫn luôn khao khát được công nhận là một thần tượng thực sự, nhưng đáng buồn thay là người hâm mộ lại chủ yếu coi cô như một ``diễn viên hài'' hay ``con thỏ điên''. Mặt khác, cô cũng đặc biệt nổi tiếng vì có thể gần như ngay lập tức chuyển đổi thái độ từ hả hê tự mãn sang tuyệt vọng cầu xin, tùy thuộc vào việc cô có chiếm ưu thế trong tình huống đó hay không.

Cô có một điệu cười đặc trưng (thường được thể hiện trong văn bản là ``AH\downarrow HA\uparrow HA\uparrow HA\uparrow'') nghe hơi giống tiếng cười của một quý tộc nữ lớn tuổi, và thường luôn thêm cụm ``-peko'' vào cuối mỗi câu. Thậm chí, điệu cười của cô còn có thể nâng lên một tầm cao áp đảo hơn: ``PE\arrow[45] KO\arrow[-45] PE\arrow[45] KO\arrow[-45] PE\arrow[45] KO\arrow[-45]\footnote{Xem phần {\abv Bạn có biết?} ở cuối tập này để đọc hướng dẫn viết điệu cười của Pekora trong \LaTeX.}.''

Tóc của Pekora dài đến giữa lưng và mang màu xanh da trời tết thành hai bím, điểm xuyết những sợi màu trắng, đặc biệt là luôn có cà rốt được cài khéo léo vào trong lọn tóc. Cô quàng một chiếc khăn lông hình con thỏ quanh cổ, vốn thực ra là một sinh vật thỏ sống có tên là Don-chan. Đôi mắt cô dường như mang hai màu tím và cam, với con ngươi hình thỏ màu trắng độc đáo. Trang phục của Pekora gồm một chiếc áo choàng (trông giống váy) viền lông màu trắng mặc bên ngoài một bộ đồ thỏ đen không dây với đường cắt dưới ngực hoặc đường cắt trên rốn, đi cùng với tay áo ngắn phồng tách rời màu trắng, găng tay viền lông đen, quần tất đen, băng nịt xếp (brial garter), băng nịt mắt cá chân bằng lông và đôi giày mary janes màu trắng có đế xanh da trời.

\chrint

``Em ngửi thấy mùi Comiket mùa đông thoang thoảng đâu đây\dots'' Korone bắt đầu hít hà đánh hơi xung quanh như một chú chó nghiệp vụ thực thụ.

Pekora ngơ ngác chớp mắt hỏi: ``Nhưng mùi ở đâu cơ chứ, peko?''

Noel hồn nhiên đáp lời: ``Mình cảm nhận nó ngay trong\dots\ cuống họng này nè.''

Lời nói đùa vẩn vơ của nàng Hiệp sĩ chưa kịp gây cười thì đã khiến cô lập tức hứng trọn một cú vung bím tóc chứa củ cà rốt đầy uy lực từ nàng Thỏ. Noel chỉ kịp xuýt xoa thốt lên hai chữ ``Ái da,'' còn cô bạn thỏ thì trừng mắt lườm, kiểu như muốn dằn mặt: ``Tôi đang cực kỳ nghiêm túc đấy nhé, peko!''

Kuon đứng cạnh chỉ biết cười khổ lắc đầu: ``Anh không nghĩ Koro-chan đang nói đùa bóng gió đâu. Có khi ẻm ngửi thấy thật theo nghĩa đen đấy.''

Nàng Khuyển vẫn miệt mài đánh hơi quanh phòng như một chú chó nghiệp vụ chuyên nghiệp, rồi bất ngờ khựng lại, chỉ thẳng ngón tay về phía bức tường kiên cố trước mặt cả nhóm: ``Ở kia kìa! Đúng chỗ đó luôn!''

\img{3-5a}

Pekora nhìn theo hướng chỉ tay, gàn giọng thắc mắc với đôi tai thỏ rung rinh bối rối: ``Đó rõ ràng chỉ là bức tường bê tông thôi mà, peko.''

Cậu Tổng nghiêm túc xoa cằm suy nghĩ, ánh mắt đăm chiêu nhìn chằm chằm vào bức tường: ``Nhưng biết đâu được\dots\ ẩn đằng sau bức tường đó lại có thứ gì đó thì sao?''

Cậu liếc sang nàng Hiệp sĩ và cất giọng đề nghị: ``Noel-chan, em thử dùng sức phá bức tường đó đi xem nào? Em khỏe nhất ở đây mà.''

Nghe yêu cầu có phần hơi quá đáng (và cũng vô cùng khác thường) từ cậu quản lý, Noel thoáng chần chừ ngần ngại: ``Hả? Anh đang nói đùa hay nghiêm túc đấy? Bắt đền em thì sao?''

Cậu đáp lại bằng một cái nhún vai tỉnh rụi: ``Yên tâm đi, có gì anh sẽ viết báo cáo đền bù thiệt hại cho công ty. Cứ mạnh dạn thử đi, để xem cái mũi của Korone-chan có nói đúng không.''

Cô thở dài cái thượt. Nhưng lời đảm bảo chắc nịch của cậu con trai dường như đã gỡ bỏ hoàn toàn sự e ngại trong lòng cô. Noel bước tới sát bức tường, cẩn thận đặt tay lên, gõ nhẹ vài cái như để thăm dò, rồi cất tiếng gọi: ``A-nô\dots\ có ai ở trong đó không vậy?''

Ngay sau câu hỏi bâng quơ đó, bức tường bỗng nhiên nứt toác rồi đổ sập xuống ầm ầm, khiến nàng Hiệp sĩ giật mình nhảy lùi lại. Lớp bụi mờ mịt dần tan đi, để lộ ra cảnh tượng Shion đang bị kẹt cứng trong bức tường với tư thế không thể nào kỳ quặc và tức cười hơn. Xung quanh cô, hàng tá những cuốn ca-ta-lô quảng bá của \hll\ cũng đang ghim chặt trên mảng tường vỡ, với vài cuốn rớt lả tả xuống sàn.

\img{3-5b}

Cả Kuon và nàng Hiệp sĩ đều trố mắt kinh ngạc, trong khi nàng Thỏ chỉ biết ngơ ngác nhìn đàn chị đầy khó hiểu. Còn nàng Chó nâu thì nhảy cẫng lên reo hò: ``Đó! Em đã bảo rồi mà! Không sai đi đâu được!''

``Tại sao chị lại bị kẹt trong tường thế này!? Rốt cuộc chuyện này là sao hả, Shion-senpai?'' Noel hốt hoảng hỏi.

``Với cả, mấy cuốn ca-ta-lô kia mắc mớ gì lại nằm ở đó?'' Kuon nhíu mày thắc mắc.

Korone nghiêng đầu nhìn đàn chị, che miệng cười khúc khích: ``Trông dáng vẻ của bà lúc này ngố tàu thật đấy á\~{}''

Nàng Phù thủy đảo mắt nhìn cả nhóm, khuôn mặt ửng đỏ vừa xấu hổ vừa cố tỏ ra bình thường thanh lịch: ``À thì\dots\ chị đang hăng say luyện tập phép dịch chuyển tức thời\dots\ nhưng mà hơi trượt tay tính toán sai tọa độ nên là\dots''

Cậu khoanh tay trước ngực, nhíu mày nghiêm nghị chất vấn: ``\dots Thế nên mới kẹt nửa người trong tường như này hả? Thế còn mấy cuốn ca-ta-lô thì sao? Anh không nghĩ chúng nó mọc cánh tự chui vào tường đâu.''

Shion cười gượng gạo, toát mồ hôi hột: ``Ừm\dots\ thì\dots\ em dùng chúng làm vật mẫu để thực hành đó. Ai mà ngờ sức mạnh của em nó lại thành công đến mức\dots\ vượt cả mong đợi như này!''

Kuon vỗ trán cái đốp, thở dài ngao ngán: ``Thành công tới mức kẹt cứng trong tường thế này\dots\ Em làm anh cạn lời luôn rồi đấy, Shion-chan ạ.''

Cậu bước tới, nắm lấy tay áo định kéo nàng phù thủy ra khỏi mớ bê tông lộn xộn, nhưng rồi khựng lại hỏi: ``Mà này, thế quái nào em có thể kẹt được với cái tư thế siêu thực như thế? Cả đầu, cả tay chân đều bị ép sát mép tường thế này á!?''

``Làm sao mà em biết được cơ chứ!'' nàng Phù thủy bật lại gắt gỏng, ``Đến em còn không hiểu nổi tại sao nó lại thành ra bộ dạng thế này nữa!''

Cả nhóm chỉ biết nhìn nhau, nửa buồn cười, nửa bối rối trước sự vụ không thể nào kỳ lạ hơn này.

``Đúng là đồ ngốc nghếch mà, peko,'' Pekora bĩu môi mỉa mai đàn chị. Đến cậu Tổng cũng phải gật gù đồng tình với cô: ``Peko-chan phán câu này chuẩn đấy. Không cần phải bàn cãi gì thêm.''

Shion ngay lập tức nhăn mặt, giọng pha chút bức xúc tự ái: ``Ý anh là sao hả, Kuon-senpai!? Mà còn em nữa nhé, Pekora-chan, em có tư cách gì mà nói chị hả? Chẳng phải bản thân em cũng đang có mấy củ cà rốt kẹt dính trong tóc đấy thôi?''

\img{3-5c}

``Hai cái đó bản chất hoàn toàn khác nhau đó em ạ,'' Kuon thở dài giải thích, ``một cái là do tai nạn phép thuật ngu ngốc, còn một cái là do sở thích cá nhân người ta tự nguyện nhét vào.''

Nàng Thỏ hừ một tiếng tự đắc, cảm thấy bị ``đụng chạm đến lòng tự tôn'', cô liền vểnh tai lên cất lời khoe khoang: ``Kuon-senpai nói chí phải lắm, peko! Không giống như chị Shion-senpai hậu đậu đâu, mấy củ cà rốt của em là tuyệt tác nghệ thuật, có thể gỡ ra gắn vào bất cứ lúc nào tùy theo ý muốn, peko!''

``Thế á? Chém gió vừa thôi. Vậy em thử rút ra đi cho bọn anh xem,'' cậu nhướng mày thách thức, tỏ vẻ không tin.

Pekora hất mặt tự tin, đưa tay lên khéo léo gỡ nhẹ củ cà rốt khỏi lọn tóc đang tết để chứng minh. Nhưng vừa lúc củ cà rốt được rút ra, một tiếng rít xé tai vang lên. Đột nhiên, một hố đen không gian mini xuất hiện ngay từ chỗ lọn tóc trống, tạo ra sức hút khổng lồ nuốt chửng mọi thứ vào trong: Noel, Korone, Shion và cả đống cuốn ca-ta-lô bay lả tả. Chỉ có Kuon là vẫn đứng vững chôn chân tại chỗ, tay ôm đầu chết sững vì bất ngờ.

\begin{figure}[H]
  \centering
  \begin{subfigure}[t]{0.45\textwidth}
    \centering
    \includegraphics[width=8cm]{../../../../Image/Book3/3-5d2.png}
    \caption{Pekora rút cà rốt ra\dots}
  \end{subfigure}
  \quad
  \begin{subfigure}[t]{0.45\textwidth}
    \centering
    \includegraphics[width=8cm]{../../../../Image/Book3/3-5d1.png}
    \caption{\dots và thảm họa bắt đầu.}
  \end{subfigure}
\end{figure}

Nàng Thỏ há hốc mồm ngạc nhiên khi chứng kiến mọi người bị lốc xoáy mini cuốn sạch sành sanh vào trong đầu mình. Cô ngó nghiêng ngó dọc: ``Ể\dots?''

Còn cậu tổng, đứng chứng kiến cảnh tượng siêu thực khó tin ấy, chỉ biết lẩm bẩm trong vô thức: ``\dots Cái\dots\ quái gì vừa xảy ra vậy?''

Cậu xoa cằm, nhíu mày nhìn chằm chằm vào khoảng không trước mặt cùng với nàng Thỏ, rồi lắc đầu cảm thán: ``\textit{Giờ thì mình đã hiểu tại sao em ấy lúc nào cũng phải kẹp khư khư củ cà rốt trên tóc rồi. Hóa ra cái gì cũng có lý do của nó cả\dots\ nhưng nút phong ấn lỗ đen kiểu này thì kỳ quặc quá.}''

\ct

Phía bên kia hố đen, Noel, Korone, và Shion rơi đánh *BIJXH* xuống một khu vực kỳ lạ. Nơi này rộng thênh thang, mọi thứ xung quanh đều mang sắc xanh lam nhàn nhạt và cấu tạo từ vô vàn những sợi khổng lồ mềm mại, như thể họ đang lạc bước trong một thế giới mơ hồ.

\img{3-5e}

Noel lồm cồm bò dậy, phủi bớt lớp bụi dính trên quần áo. Cô ngó nghiêng xung quanh không gian kỳ ảo này, giọng đầy lo lắng: ``Chúng ta\dots\ rốt cuộc đang ở đâu thế này?''

Korone lập tức nhảy lên vài cái, hít hà không khí xung quanh như một chú chó nghiệp vụ đang đánh hơi tìm đường: ``Không biết nữa, nhưng chị cảm nhận được cái mùi gì đó\dots\ rất thân quen. Giống y hệt mùi của Pekora-chan!''

Trong khi đó, Shion đang dùng ma thuật để lơ lửng giữa không trung quan sát địa hình. Cô nhíu mày, đưa ra một giả thuyết đầy tính logic phép thuật: ``Chắc chắn là chúng ta đã bị hút vào cái lỗ hổng không gian bên trong lọn tóc của Pekora để lấp chỗ trống cho củ cà rốt vừa bị rút ra rồi.''

Bất ngờ, một giọng nói quen thuộc vang vọng từ hư không vọng tới: ``Ừm, phân tích nghe cũng có lý lắm.''

Nàng Phù thủy giật bắn mình, quay phắt lại, và\dots\ đối diện trực diện với một con mắt đỏ khổng lồ vĩ đại đang nhìn chằm chằm vào mình xuyên qua bầu trời.

Hoảng hốt tột độ, cô mất tập trung ma thuật và suýt rơi cắm đầu xuống đất nếu không kịp quẫy đạp lấy lại thế thăng bằng. Nhận ra chủ nhân của con mắt to tổ chảng đó là ai, cô liền bực tức hét lớn vang trời: ``\textbf{M-Mồ! Kuon-senpai! Đừng có xuất hiện rình rập làm em giật mình chết khiếp như thế chứ!!}''

Kuon, lúc này từ góc nhìn của họ chỉ là một con mắt to tròn, cười gượng gạo vang vọng: ``À, anh xin lỗi, anh không cố ý dọa ma em đâu.''

Noel, vẫn còn ôm ngực thót tim sau cú giật mình bởi con mắt to đùng kia, liền tò mò hỏi vọng lên: ``Mà, sao anh biết bọn em đang ở đây vậy, Kuon-senpai?''

Con mắt vĩ đại chớp chớp đáp lại: ``Khó gì đâu. Cứ ghé mắt dòm vào đúng cái lỗ mà Pekora-chan vừa rút củ cà rốt ra là thấy rõ mồn một ngay. Thực ra mà nói, mấy đứa bị thu nhỏ đi trông bé xíu như đàn kiến đang chạy rành rành luôn.''

Cô Hiệp sĩ vừa nghe vừa gật gù, nhưng ánh mắt không giấu nổi sự thắc mắc: ``Chả trách sao trông anh to khủng bố thế này\dots\ Nhưng mà, anh ở tít bên ngoài làm sao mà nghe rõ bọn em nói nhỏ xíu thế này?'

Cậu Tổng cười đáp lại một cách vô cùng đơn giản và thực dụng: ``À, anh đang dùng tai nghe khuếch đại âm thanh siêu việt. Đồ xịn do anh tự chế đấy, bắt sóng tốt lắm.''

Noel thở dài cái thượt, gạt nốt mớ bụi bặm vương trên áo mình, rồi ngẩng đầu lên nhìn thẳng vào con mắt khổng lồ đang lơ lửng trên đỉnh đầu: ``Thế\dots\ bây giờ bọn em phải làm gì tiếp theo đây? Anh có cách nào gắp hay kéo bọn em ra ngoài không?''

Cậu Tổng khẽ di chuyển con ngươi đảo quanh lọn tóc, như thể đang quan sát bao quát toàn bộ thế giới nhỏ bé bên trong đó. Sau một hồi suy nghĩ nát óc, cậu khẽ nhắm mắt lại, rồi mở ra nói với vẻ vô cùng bất lực: ``Cái đó thì anh chịu thua. Anh nghĩ mấy đứa nên tự thân vận động, thử đi bộ nhìn xung quanh xem sao. Biết đâu lại có cái lối thoát bí mật nào đó.''

Ở bên ngoài, Pekora đang đứng yên như tượng với vẻ mặt vừa bực bội vừa sốt ruột, liền lên tiếng gắt gỏng: ``Nè, giải quyết xong chưa đó? Em chịu hết nổi rồi nha, peko! Mỏi chân quá đi mất!''

Kuon, giờ đang khom người đứng sát bên cạnh nàng Thỏ, cứ cúi gập đầu ghé mắt nhìn chằm chằm sát rạt vào lọn tóc của cô. Còn bản thân nàng Thỏ thì đang phải kiễng chân khổ sở, cố gắng giữ thăng bằng không dám nhúc nhích, hai tay chống hông đầy vẻ khó chịu.

Cô cằn nhằn cất giọng: ``Kuon-senpai! Anh mau dùng não bộ quản lý nghĩ cách đi chứ! Chẳng lẽ anh định bắt em đứng làm tượng sáp vĩnh viễn ở đây hả, peko?''

Cậu Tổng đáp lại tỉnh bơ, mắt vẫn không rời khỏi lọn tóc tí ti: ``Bình tĩnh đi em, đừng có nóng nảy. Anh đang vận nội công suy nghĩ làm việc đây chứ gì nữa.''

Nàng Thỏ lầm bầm chửi rủa gì đó không rõ trong miệng, nhưng vẫn cam chịu cố giữ nguyên tư thế bất động.

\ct

Bên trong thế giới lọn tóc, ba người họ tiếp tục đi lang thang và suy nghĩ cách thoát ra khỏi chốn giam cầm này.

Shion bay lơ lửng, cất giọng phân tích tình hình một cách vô cùng bình tĩnh: ``Để xem nào\dots\ Có lẽ chúng ta nên thử dùng ma thuật tìm một điểm kết nối không gian ở đâu đó, hoặc ít nhất là--''

Cô nàng tóc xám bỗng dưng nghẹn họng dừng lại giữa câu khi nhận ra hai người đồng hành của mình, Korone và Noel, từ nãy giờ chẳng thèm lọt tai chút xíu nào những gì cô đang nói. Cả hai đang chăm chú\dots\ ngồi bệt xuống vuốt ve thảm tóc của Pekora một cách vô cùng đắm say và thích thú.

\img{3-5f}

Nàng Hiệp sĩ, tay vuốt nhẹ từng sợi tóc mềm mại khổng lồ như vuốt thảm lông cừu, không kìm được mà trầm trồ khen ngợi: ``Ôi chao mượt quá đi mất\~{} Đúng là chất tóc được chăm sóc kỹ lưỡng của Peko-chan có khác! Nằm sướng rơn người!''

Nàng Khuyển mắt sáng rực rỡ, reo lên phấn khích: ``Nhìn sợi tóc này bóng bẩy chưa kìa! Trông ngon miệng cứ y như món cơm chan trà xanh ochazuke vậy đó!''

Nàng Phù thủy trừng mắt, cảm giác như mình vừa bị ``đá văng'' không thương tiếc khỏi cuộc họp bàn chiến lược sinh tử. Cô cố gắng quơ tay quơ chân múa may để thu hút sự chú ý của hai con người đang tận hưởng kia: ``Làm ơn tập trung nghe chị nói đi mà! Chúng ta đang mắc kẹt nguy hiểm trong đây đấy, đừng có mà đùa giỡn cợt nhả nữa!''

Nhắc đến món ochazuke và thức ăn, ánh mắt Noel bỗng chốc trở nên lờ đờ mơ màng. Một tay cô ôm chặt lấy cái bụng đang sôi réo rắt, vẻ mặt nhăn nhó khổ sở: ``Nhắc mới nhớ\dots\ Mình đói rã ruột quá rồi\dots''

Nàng Chó lạc quan nhanh chóng đề xuất ý tưởng: ``Thử chia nhau tìm quanh khu vực này xem! Biết đâu em ấy lại giấu củ cà rốt dự trữ nào đó chôn sâu trong mớ tóc này thì sao?''

Dù ý tưởng đó nghe qua chẳng mấy hợp lý hay thuyết phục tẹo nào, nhưng với một tâm hồn ăn uống như Noel, chỉ cần nhắc đến hai chữ đồ ăn cũng đủ làm cô sáng bừng sức sống. Không cần suy nghĩ do dự thêm giây nào, cô hớn hở bật dậy như lò xo: ``Quyết định vậy đi! Để em xông pha đi tìm thử! Cà rốt hay củ cải rễ cây gì cũng nhai được hết!''

Vừa hừng hực khí thế xoay người quay đi chưa được vài bước, Noel lại bất cẩn vấp chân phải một lọn tóc tết xoắn ốc thò ra. Cô loạng choạng mất đà rồi ngã chúi nhủi mặt đập xuống sàn, khiến Korone đứng phía sau không những không đỡ mà còn ôm bụng bật cười khoái chí: ``A là la là la\~{} Cú ngã 10 điểm đẹp mắt lắm, Noel-chan!''

\img{3-5g}

Những tưởng cô sẽ bị u đầu dập mũi và nằm bẹp dí trên nền tóc mềm xèo, nhưng kỳ lạ thay, toàn thân cơ thể cô bỗng lóe lên luồng ánh sáng chói lóa rồi tan biến vào hư không.

Trong cái chớp mắt ngắn ngủi, cô đã bị đẩy văng từ không gian bên trong lọn tóc của nàng Thỏ tuột ngược ra ngoài, rơi phịch xuống sàn nhà trở lại vị trí ban đầu của căn phòng nghỉ, ngay trước ánh mắt mở to ngạc nhiên tột độ của cả Pekora và Kuon.

Cậu Tổng thấy Noel đột nhiên xuất hiện ``từ trên trời rơi xuống'' thì vội vàng bước tới đỡ, ngạc nhiên hỏi dồn dập: ``Làm thế nào mà em tự phá được kết giới thoát ra ngoài nhanh thế?''

Nàng Hiệp sĩ phủi phủi váy ngồi dậy, vẻ mặt đờ đẫn ngơ ngác nhìn xung quanh, như thể chính cô cũng không hiểu phép màu nào đã đưa mình trở lại thế giới thực: ``Em cũng chả biết nữa anh ơi\dots\ Em chỉ định chạy đi kiếm đồ ăn, rồi tự nhiên vấp ngã sấp mặt\dots\ Rồi mở mắt ra tự dưng lại thấy nằm ở đây.''

Cậu xoa cằm đưa tay lên gãi đầu khó hiểu: ``Thế thì\dots\ Anh cũng cạn lời, không biết phải giải thích theo tính vật lý hay ma thuật nữa.''

Cô lảo đảo chập choạng bước đi được vài bước rồi lả dần đi vì cơn đói cồn cào ập đến. Cô ôm khư khư cái bụng than vãn thảm thiết: ``Kuon-senpai ơi, ai đó làm ơn nấu cái gì đó cho em bỏ bụng đi\dots\ Em đói sắp chết ngất đến nơi rồi.''

Kuon nhún vai bất đắc dĩ, xắn tay áo lên rồi quay sang Pekora đang đứng yên như tượng: ``Thôi được rồi, để anh xuống bếp nấu bát mì cứu đói cho mọi người ha. Pekora-chan có muốn ăn một bát không?''

Nàng Thỏ nuốt nước bọt, gật đầu lia lịa như gà mổ thóc, dù vẫn cố ra vẻ điềm tĩnh thanh cao: ``Ăn chứ, có ngu mới không ăn, peko! Anh định nấu mì ramen với topping gì vậy, peko?''

Cậu Tổng liệt kê thực đơn: ``Mấy cái nguyên liệu lặt vặt còn thừa trong tủ lạnh thôi. Vài quả trứng, chút thịt nạc nếu em ăn được, và ít rau củ, đặc biệt là \emph{cà rốt} thái sợi gì đấy.''

Lập tức, đôi tai thỏ của cô vểnh đứng lên rung rinh đầy phấn khích: ``Nghe duyệt đấy! Vậy thì anh mau làm nhanh lên đi, peko!''

\ct

Sau khi nấu xong bát mì nóng hổi thơm phức dâng tận miệng cho hai cô gái đang lả đi vì đói, Kuon lại tiếp tục nhiệm vụ ghé sát mắt ngó vào trong cái lỗ lọn tóc của Pekora để hướng dẫn Korone và Shion tìm đường trở về. Bên trong thế giới thu nhỏ ấy, hai người họ đang bấn loạn vắt óc tìm đủ mọi cách để giải cứu Noel, trong khi chính người cần cứu thì đang thong thả ngồi ngoài húp mì xì xụp.

Nàng Phù thủy vẫn chưa hết bàng hoàng ngơ ngác sau sự biến mất chớp nhoáng đầy bí ẩn của nàng Hiệp sĩ, nhưng cô cố hít một hơi sâu để trấn tĩnh lại tinh thần: ``Khoan đã nào\dots\ rốt cuộc thì sao mà em ấy bị bay màu quay về được thế này!? Cơ chế ở đây là gì chứ!?''

Cậu Tổng, giờ đây đang bưng nguyên bát mì nóng hổi trên tay, vừa húp xì xụp vừa nói vọng vào trong: ``Theo giả thuyết của anh, chắc tại lúc vấp ngã thì Noel-chan đã vô tình chạm phải cái cuốn ca-ta-lô. Cuốn sách đó thuộc về không gian bên ngoài nên nó có tính năng đẩy người ta ra. Anh đoán bừa thế.''

Nàng Khuyển gật gù ra vẻ đắc ý như đã thấu hiểu chân lý vạn vật. Cô quay sang nhìn Shion, hỏi: ``Chị có đem theo món đồ ma thuật nào có khả năng tương tự không?''

Nàng Phù thủy lôi từ trong áo choàng ra mấy tờ giấy, khai nhận với vẻ thất vọng: ``Lúc nãy đi vội, trên người chị chỉ có cầm dư mấy cuốn ca-ta-lô quảng bá cho Comiket mùa đông thôi.'' Nói đoạn cô giơ chúng lên cho hai người xem.

Lập tức, Korone búng tay cái chóc, reo lên đắc thắng: ``Tuyệt vời! Vậy thì dùng luôn đống sách đó làm chìa khóa đi!''

Cô nhào tới giật phăng lấy một cuốn ca-ta-lô trên tay của nàng ``Tỏi'', rồi dậm chân nhảy lên ấn mạnh cuốn sách cắm phập xuống lớp sàn tóc. Ngay tắp lự, không gian xung quanh rung chuyển. Cuốn ca-ta-lô ấy lóe sáng rồi trồi ngược lên, bay vọt ra khỏi lọn tóc của nàng Thỏ trước sự kinh ngạc tột độ của cả ba người.

\begin{figure}[H]
  \centering
  \begin{subfigure}[t]{0.45\textwidth}
    \centering
    \includegraphics[width=8cm]{../../../../Image/Book3/3-5h1.png}
    \caption{Shion đưa ra hai cuốn ca-ta-lô, những gì còn lại trên người\dots}
  \end{subfigure}
  \quad
  \begin{subfigure}[t]{0.45\textwidth}
    \centering
    \includegraphics[width=8cm]{../../../../Image/Book3/3-5h2.png}
    \caption{\dots Korone lấy một trong hai cuốn đó ấn xuống\dots}
  \end{subfigure}
  \quad
  \begin{subfigure}[t]{0.45\textwidth}
    \centering
    \includegraphics[width=8cm]{../../../../Image/Book3/3-5h3.png}
    \caption{\dots để rồi chính nó quay trở về văn phòng.}
  \end{subfigure}
\end{figure}

Ở thế giới thực, Pekora đang một tay cầm thìa chuẩn bị rót trà, còn Noel thì đang gắp cuộn mì thơm phức mọng nước do cậu nấu đưa lên miệng. Cả hai người bọn họ khựng lại, tròn mắt nhìn cuốn ca-ta-lô từ đâu tự dưng phóng vụt ra, mắc kẹt lủng lẳng trên lọn tóc của nàng Thỏ. Họ chỉ biết đưa mắt nhìn nhau, đồng thanh bật ra một tiếng ``Ể?'' đầy mông lung.

Kuon đứng bên cạnh, vẫn đang ung dung xì xụp húp nước dùng bát mì, nhai nhóp nhép khen ngợi: ``Anh nghĩ cách kết nối vật lý này có vẻ hiệu quả thực sự đấy! Hai đứa ở trong cứ áp dụng chiêu này rồi dùng pháp thuật kết nối nó với cánh cổng gì đó xem sao!''

Bất chợt, nghe cậu Tổng nhắc đến pháp thuật, nàng Chó bỗng gõ đầu cái cốc nhớ ra một điều vô cùng quan trọng. Cô vội vã quay ngoắt sang túm vai nàng Phù thủy, hai mắt sáng rỡ cất giọng hỏi: ``Nghĩ lại mới nhớ nha, Shion-chan, ngay từ đầu cái tai nạn sập tường là do chị đang rèn luyện phép dịch chuyển không gian cơ mà đúng không?''

Shion khựng lại ngẫm nghĩ một lúc, rồi cũng vỡ lẽ sực nhớ ra mục đích ban đầu định làm từ trước khi bị kẹt dính trên tường. Cô chống hông tự tin, vênh mặt hậm hực đáp lại: ``Phải rồi ha! Được thôi, để pháp sư vĩ đại đây ra tay thử xem. Nhưng nói trước nha, nếu lần này lỡ trượt tay thất bại nữa làm chúng ta bay ra ngoài vũ trụ thì chị đây không chịu trách nhiệm đâu đấy!''

\img{3-5i}

Dứt lời, nàng Phù thủy bắt đầu tập trung cao độ niệm chú vận phép. Từ dưới chân họ, một vòng tròn ma pháp trận màu tím khổng lồ rực rỡ hiện ra. Cả khu vực không gian bên trong lọn tóc đều lóe sáng rực lên lung linh chói lóa.

Nàng Chó lại nhảy múa reo hò phấn khích không ngừng: ``A là la là la\~{} Phép thuật của Shion-chan trông lấp lánh thật là đẹp quá xá luôn!''

Cột sáng phép thuật của cô rọi lên sáng rực, sáng mãi, để rồi một tiếng *PỤP* vang lên\dots

\ct

\dots Cả cô và Korone đều được dịch chuyển bình an vô sự, hạ cánh an toàn quay trở về căn phòng nghỉ thân thương, trước sự vỗ tay vui mừng khôn xiết của Noel.

Kuon bưng bát mì trống rỗng, thở phào nhẹ nhõm một hơi dài như trút được gánh nặng vì cậu đinh ninh rằng mọi rắc rối cuối cùng cũng đã êm xuôi kết thúc: ``\textit{Phước đức quá đi mất. Cũng may là trình độ phép thuật cùi bắp của em ấy lần này lại có hiệu quả vào đúng lúc then chốt.}''

\img{3-5j}

Nàng Chó mừng rỡ chạy ùa tới ôm chầm lấy nàng Hiệp sĩ, hai người nắm tay nhau xoay vòng nhảy vũ điệu chiến thắng điên cuồng. Nàng Phù thủy vuốt lại mái tóc, đứng nhìn hai cô bạn vui vẻ ăn mừng, trên môi cũng nở một nụ cười mãn nguyện đầy tự hào như muốn chúc phúc cho chiến tích lẫy lừng của họ.

``Khoan đã, từ từ\dots'' bất chợt, cậu Tổng cau mày cất giọng cắt ngang bầu không khí ăn mừng. ``Pekora-chan đâu rồi?''

Nghe vậy, cả nhóm giật thót mình, vội vàng xúm lại đảo mắt tìm kiếm xung quanh. Cô Thỏ tóc xanh không hề ngồi ngoan ngoãn trên ghế ăn mì cùng với Noel như trước đó, cũng lặn mất tăm khỏi khu vực trung tâm mà họ đang đứng tụ tập.

``Lạ thật đấy, rõ ràng là em nhớ in dấu não là cậu ấy vừa nãy đứng ở ngay đây với em và Kuon-senpai ăn mì mà,'' Noel ngơ ngác gãi đầu.

``Ừ nhỉ, kỳ cục thật. Chị cứ đinh ninh tưởng là khi nãy thi triển pháp thuật thì em ấy cũng phải dịch chuyển về đứng ở đây đón chúng ta chứ?'' Shion nhíu mày thắc mắc.

Korone không nói lời nào, chỉ rướn cổ khịt khịt mũi đánh hơi không khí. Nàng Khuyển tiếp tục kích hoạt tối đa khứu giác siêu việt của mình để dò tìm mục tiêu. Ngay sau đó chưa đầy ba giây, cô giơ tay reo lên: ``Có mùi nồng nặc ở đằng kia kìa!'' và chỉ thẳng ngón tay về phía bức tường nứt nẻ, nơi mà cách đây không lâu chính Shion vừa bị mắc kẹt chui ra.

Tất cả mọi người từ từ quay đầu nhìn sang, đứng ngớ người ra một lúc, rồi đồng loạt ôm bụng bật cười sằng sặc không phanh trước tình cảnh thê thảm hiện tại của Pekora.

Toàn thân cô Thỏ tội nghiệp giờ đây đang bị ép chặt cứng ngắc trên bức tường, xung quanh cắm lỉa chỉa đầy mấy củ cà rốt quen thuộc. Tay cô khư khư cầm chặt một cuốn ca-ta-lô Comiket làm lá chắn, trong khi đôi môi mím chặt, cả thân hình thỏ cứ run lên bần bật vì vừa tức vừa sợ. Vòng tuần hoàn kẹt tường lại một lần nữa tiếp diễn một cách hoàn hảo!

\img{3-5k}

%==========================================TRUYỆN PHỤ=========================================
\omk

Khi chứng kiến nàng Thỏ kiêu hãnh đang bị dính chặt trong tư thế ``khó đỡ'' vô cùng thảm hại trên tường, cả bốn người còn lại thực sự không thể nào nhịn nổi tràng cười đứt ruột.

Pekora, mặt đỏ phừng phừng lên đến tận mang tai vì quá sức xấu hổ, gân cổ gào thét đến lạc cả giọng: ``\textbf{Mấy anh chị cười cái khỉ khô gì chứ!? Cười đủ chưa hả? Mau tới đây lôi em thoát khỏi cái chỗ quỷ quái này đi, peko!}''

Shion, người là nạn nhân tiền bối từng nếm trải cảm giác này nên có vẻ đồng cảm và cười nhẹ nhàng nhất đám, cố gắng bịt miệng trấn tĩnh lại và lau nước mắt nói lời xin lỗi: ``Ừm\dots\ chị xin lỗi nha Peko-chan\dots\ Nhưng mà, thú thật là nhìn bộ dạng em thế này buồn cười quá mức quy định rồi!''

Trong khi đó, Kuon đứng khoanh tay, tranh thủ lúc cô nàng đang kẹt cứng để mở lớp giáo huấn đạo lý trong lúc bản thân cũng đang cố nhịn cười vất vả: ``Thấy hậu quả chưa Shion-chan? Anh nghĩ là về sau em nên bỏ thời gian tập luyện lại kiểm soát quỹ đạo phép thuật dịch chuyển đi. Không thì sau này lại cứ đi rải rác tạo thêm mấy vụ kẹt tường như thế này nữa, mọi người lại có dịp cười bò ra nhà mất.''

Pekora nghe cậu Tổng bình phẩm vậy thì càng tức tối lộn ruột hơn, cô giãy nảy cố giằng thân mình khỏi mảng bê tông: ``Ê này, đừng có mà đổi chủ đề đổ lỗi cho chị ấy nha! Rõ ràng là lỗi của mấy người rủ rê chơi bời đã vung tay vô tình đẩy em vào cái tình cảnh thảm thiết này, peko! Đền bù thiệt hại tinh thần đi!''

Lúc này, con mắt tinh tường của đói khát của Noel mới lia tới và để ý đến mấy củ cà rốt đang đính chặt chẽ trên tường ngay sát bên cạnh chỗ Pekora. Nàng Hiệp sĩ liếm môi háo hức tiến tới gần hỏi xin xỏ: ``Nè Pekora-chan ơi, dù sao cũng kẹt rồi, cho bọn mình xin chia lại mấy củ cà rốt ngon lành này để nấu súp ăn nhé?''

Lập tức, nàng Thỏ như bị đụng chạm tới kho báu sinh mệnh, đã hét lên phản đối kịch liệt văng cả nước bọt: ``\textbf{Không đời nào! Cà rốt bảo bối của tôi mà! Cấm ai đụng vào!}''

Korone nháy mắt đầy ẩn ý, thản nhiên hùa theo chọc ghẹo thêm cho vui: ``Trời ơi ki bo quá nha\~{} Pekora-chan tích trữ giấu giếm được nhiều cà rốt thế mà không chịu chia cho đồng đội đang đói khát gì cả. Xấu tính ghê nơi! Em là em chia đều cho mọi người rồi đó!''

Đến cả cậu Tổng lúc này cũng không tha, hăng hái đổ thêm dầu vào ngọn lửa đang bùng cháy: ``Đúng đó nha! Tinh thần đồng đội ở đâu hả? Có đồ ngon thì phải biết chia sẻ bát cơm manh áo với mọi người chứ, Peko-chan! Em giấu ăn một mình kẻo lại mập thỏ đó!''

Nàng Thỏ lúc này bị khích tướng đâm ra càng tức tối lộn tiết hơn, giọng gân lên the thé như muốn cãi tay đôi sống mái với cả thế giới: ``\textbf{Mấy người thì biết cái gì cơ chứ! Có hiểu là bộ tộc loài thỏ chúng tôi sinh ra và sống là vì cà rốt không hả? Cà rốt chính là cội nguồn sức mạnh và lẽ sống của tôi mà, peko!}''

Cậu quản lý nghe luận điệu phi thực tế ấy thì nhướng một bên mày, lập tức giở giọng mọt sách kiến thức ra phản bác ngay tắp lự: ``Xin lỗi cô nương, bớt xạo lại đi. Thỏ ngoài đời thật thực chất không có phát cuồng vì cà rốt đến mức thần thánh hóa vậy đâu. Dạ dày chúng yếu lắm, thường chúng chỉ khoái ăn rau xanh với cỏ khô thôi. Rốt cuộc em là cái thể loại thỏ lai giống gì mà lập dị khác đời thế?''

Bị bóc mẽ sự thật phũ phàng, cô nàng nhất quyết không chịu thua, tức điên máu lên gào thét: ``\textbf{Rõ ràng trong anime là thế mà! Mà thôi dẹp đi! Rốt cuộc thế mấy người tóm lại có định giơ tay giúp lôi tôi ra không, hay là rảnh rỗi chỉ đứng đó trêu chọc bóc phốt tôi hả, peko!?}''

Cậu cười khẩy một tiếng nham hiểm, khoanh tay đứng nhìn từ trên xuống: ``Giúp thì cũng được thôi, chuyện nhỏ. Nhưng quy luật có vay có trả, ít nhất em phải nộp lệ phí chia đôi số cà rốt kia cho bọn anh làm lẩu chứ nhỉ? Không thì cứ từ từ mà tận hưởng phong cảnh vách tường đi.''

Pekora trừng mắt, nghiến răng đáp lại cậu bằng một giọng nói đanh thép không lay chuyển: ``Kẻ sĩ thà chết vinh còn hơn sống nhục! Có kẹt chết khô thành xác ướp ở đây thì tôi cũng không chia một mẩu, peko!''

Thế rồi, nhận thấy cãi lý không lại, cô nàng quyết định giở bài chuồn, bắt đầu ``bật chế độ'' ăn vạ đỉnh cao. Cô rên rỉ liên hồi như con nít đòi quà, khóc lóc ỉ ôi nước mắt ngắn nước mắt dài làm nũng ỏm tỏi, khiến cả nhóm chẳng những không động lòng mà còn không thể nhịn nổi buồn cười. Họ thay nhau buông lời trêu chọc, chọc tức cô tới bến nhiều hơn nữa.

Cuối cùng, gồng mình chịu trận không nổi trước sức ép đàm tiếu từ mọi phía, Pekora đành ngậm đắng nuốt cay miễn cưỡng nhượng bộ, thút thít hứa chia cho mỗi người hai củ cà rốt tươi ngon. Đây chính là ``cái giá đắt đỏ'' mà cô phải trả để nhóm người này đồng ý vươn tay giải thoát cô ra khỏi bức tường oan trái.

Sau khi đã trêu đùa thỏa thích và thỏa thuận xong xuôi hợp đồng chia chác cà rốt, Kuon lấy điện thoại ra, quyết định gọi dịch vụ sửa chữa tường chuyên nghiệp của tòa nhà lên xử lý. Cậu cất giọng giải thích, cố ý nhấn mạnh từng chữ để khịa: ``Anh chọn thợ chuyên nghiệp vì anh hoàn toàn không muốn nhờ cậy vào bàn tay vàng của Noel-chan đâu. Với lực đập phá bằng chùy của em ấy á, có khi giải cứu xong cái văn phòng này cũng sụp thành đống đổ nát hoang tàn mất.''

Korone và Shion lại được dịp ôm bụng cười sằng sặc trước câu nói thâm thúy của cậu. Trong khi đó, nàng Thỏ Pekora thì vẫn ngồi xụ mặt giận dỗi một góc tiếc nuối số cà rốt sắp mất, còn nàng Hiệp sĩ Noel thì phồng hai má lên phụng phịu bất mãn vì bị đánh giá thấp khả năng khống chế sức mạnh.

``Đừng có đồn bậy nha!'' Noel chu môi phản đối, ``Em không có hậu đậu vụng về đến mức đập nát cái văn phòng đâu\~{}!''

Pekora thì sụt sùi lau nước mắt cá sấu: ``*hức* Đám người mấy người đúng là một lũ xấu tính tồi tệ nhất trần đời mà, peko\dots''

Đúng là một cách mở màn cho một ngày mới ngập tràn rắc rối cười ra nước mắt, và đó cũng là cách tuyệt vời để cả nhóm \hll\ khép lại năm Ất Hợi dương lịch 2019 với đầy ắp những câu chuyện hài hước và\dots\ vô cùng hỗn loạn!

\trva

Usada Pekora có hai điệu cười đặc trưng làm nên thương hiệu của cô nàng: ``AH\downarrow HA\uparrow HA\uparrow HA\uparrow'' và ``PE\arrow[45] KO\arrow[-45] PE\arrow[45] KO\arrow[-45] PE\arrow[45] KO\arrow[-45]''.

Trong hệ thống soạn thảo \LaTeX, điệu cười kiểu 1 không yêu cầu người dùng cài đặt gói bổ trợ. Để viết điệu cười của Pekora theo kiểu 1, ta đơn giản chỉ cần viết đoạn mã:

\begin{verbatim}
    AH\downarrow HA\uparrow HA\uparrow HA\uparrow
\end{verbatim}

Với kiểu cười thứ 2 mang tính phức tạp hơn, người dùng bắt buộc cần định nghĩa một hàm \verb|\arrow| riêng biệt trong phần preamble bằng gói \textit{tikz}:
\begin{verbatim}
    \newcommand{\arrow}[1][-45]
      {\begin{tikzpicture}[baseline = -0.5ex]
          \node[inner sep=0pt,outer sep=0pt,rotate = #1] 
          (a) at (0,0)  {\rightarrow};
       \end{tikzpicture}}
\end{verbatim}

Sau đó, trong phần thân tài liệu, ta gọi hàm đó để viết ra như sau:
\begin{verbatim}
    PE\arrow[45] KO\arrow[-45] PE\arrow[45] 
    KO\arrow[-45] PE\arrow[45] KO\arrow[-45]
\end{verbatim}

\end{document}
